\documentclass[a4paper,10pt]{amsart}

\pdfinfoomitdate=1
\pdftrailerid{}
\pdfsuppressptexinfo=15

\usepackage[T1]{fontenc}
\usepackage{lmodern}
\usepackage[margin=25mm]{geometry}
\usepackage{amsmath,amssymb,mathtools,xcolor}
\usepackage{array}
\usepackage[protrusion=true,expansion=false]{microtype}
\usepackage[shortlabels]{enumitem}
\usepackage[numbers,sort&compress]{natbib}
\usepackage[colorlinks=true,linkcolor=blue!55!black,citecolor=blue!55!black,urlcolor=blue!55!black]{hyperref}
\hypersetup{hypertexnames=false,pdftitle={},pdfauthor={},pdfsubject={},pdfkeywords={},pdfcreator={},pdfproducer={}}

\newtheorem{theorem}{Theorem}[section]
\newtheorem{lemma}[theorem]{Lemma}
\newtheorem{corollary}[theorem]{Corollary}
\theoremstyle{remark}
\newtheorem{remark}[theorem]{Remark}

\newcommand{\Z}{\mathbb Z}
\newcommand{\F}{\{0,1\}}
\newcommand{\wt}{\operatorname{wt}}
\newcommand{\Fix}{\operatorname{Fix}}
\newcommand{\Imap}{\operatorname{Im}}
\newcommand{\depth}{\operatorname{depth}}
\newcommand{\1}{\mathbf 1}

\title[First-frequency rotation]{First-Frequency Rotation on Binary Pointed Necklaces}
\author{Anonymous}
\date{}

\begin{document}

\begin{abstract}
Let $R$ be left rotation of a binary word and let
$m_a(w)=|\{i:w_i=a\}|$.  We study the finite map
$T_n(w)=R^{m_{w_0}(w)}w$ on $\{0,1\}^n$.  On a cyclic word of least
period $d$ and weight $k$, the pointed dynamics splits into
$\gcd(k,d)$ disjoint $\mathord{\pm}k$ generator cycles.  A constant
component is periodic; every nonconstant component flows into its directed
$10$ edges.  This gives the exact possible periods, the sharp maximum
preperiod $n-2$, and exactly two deepest states for $n\geq3$.  Independently,
we give the complete predecessor list of every target, the global
$0/1/2$ fibre distribution, and a primitive-block M\"obius fixed census.
Hamming-weight-controlled frozen rotations and the internal P166 cyclic
phase architecture receive zero contribution credit.  The retained status
is \texttt{AMBER\_INTERNAL\_NEAR\_P166 / HOLD\_EXTERNAL}.
\end{abstract}

\maketitle

\section{Literal map and subtraction boundary}

For $n\geq1$, write a word as $w=w_0w_1\cdots w_{n-1}$, with indices in
$\Z/n\Z$, and set
\[
 R(w_0w_1\cdots w_{n-1})=w_1\cdots w_{n-1}w_0.
\]
For $a\in\F$, let $m_a(w)=|\{i:w_i=a\}|$.  Our literal map is
\begin{equation}\label{eq:literal}
 T_n(w)=R^{m_{w_0}(w)}w,\qquad w\in\F^n.
\end{equation}
It preserves both the cyclic necklace and the weight $k=\wt(w)=m_1(w)$.
On a fixed weight layer its two branches are
\begin{equation}\label{eq:branches}
 w_0=1:\quad T_n(w)=R^k w,
 \qquad
 w_0=0:\quad T_n(w)=R^{n-k}w=R^{-k}w.
\end{equation}

H\o yer and \v Spalek explicitly construct the quantum phase rotation
$R_z(\varphi|x|)$ under the heading ``rotation by Hamming weight''
\cite[Sec.~3.2]{HoyerSpalek2005}.  Their object is not a cyclic coordinate
shift.  Nevertheless, we assign the Hamming-weight-controlled rotation
idea and both frozen branches in \eqref{eq:branches} zero contribution
credit.  The adaptive first-symbol gluing is the only literal feature under
analysis.

Gro\v sek--Hromada study coordinate-rotation classes at fixed binary weight,
and Gupta et al. study literal circular shifts, rotation equivalence,
weight, and complementation \cite{GrosekHromada2016,GuptaEtAl2022}.  Thus
coordinate-necklace structure, fixed-weight class enumeration, and ordinary
shift symmetries also receive zero contribution credit.  Neither source
defines the adaptive first-symbol gluing or its functional graph.

There is a closer internal boundary.  P166 studies
$S_n(x)=x+\wt(x)\1$ on $(\Z/n\Z)^n$ and, on a diagonal-translation orbit,
reduces it to $j\mapsto j+c_j$, where $(c_j)$ is a weak composition of $n$.
Our reduction below also has cyclic phase syntax, but its increments are the
two values $\mathord{\pm}k$.  The exact separation is summarized here.

\begin{center}
\small
\renewcommand{\arraystretch}{1.08}
\begin{tabular}{@{}>{\raggedright\arraybackslash}p{0.19\linewidth}>{\raggedright\arraybackslash}p{0.36\linewidth}>{\raggedright\arraybackslash}p{0.36\linewidth}@{}}
\hline
axis & map \eqref{eq:literal} & internal P166 map \\
\hline
invariant orbit & coordinate-rotation necklace & diagonal alphabet-translation orbit \\
phase profile & $a_j\in\{k,-k\}$ & $c_j\geq0$, $\sum_jc_j=n$ \\
forward engine & disjoint $\mathord{\pm}k$ generator cycles & cycle mass exhaustion \\
recurrence & possibly several nontrivial components & at most one nontrivial cycle per orbit \\
periods & $1,2$, and proper divisors of $n$ & every value $1,\ldots,n$ \\
one-step fibres & two labelled rotations; size $0,1,2$ & histogram-selected translates; unbounded in $n$ \\
\hline
\end{tabular}
\end{center}

The common phase syntax, generic functional-graph language, the numerical
clock value $n-2$, and indicator-style inverse notation all receive zero
contribution credit.  We retain only the constrained generator-component
theorem and its consequences proved below.  A direct owner of the adaptive
literal map, a literal conjugacy into P166, or a proof that the component
theorem is a formal consequence of P166 mass exhaustion triggers
\texttt{KILL\_INTERNAL\_P166\_PHASE\_MAP}.  Until that gate is resolved, the
status remains
\[
 \texttt{AMBER\_INTERNAL\_NEAR\_P166 / HOLD\_EXTERNAL}.
\]

For a point $w$ of a finite self-map, its \emph{preperiod} is the least
$t\geq0$ for which $T_n^t(w)$ is periodic.  Its eventual exact period is
called simply its period.

\section{Pointed-necklace component theorem}

Let $u\in\F^n$ have least rotational period $d\mid n$.  Its distinct
pointed states are $R^ju$, $j\in\Z/d\Z$.  We use the periodic indexing
$u_{j+d}=u_j$ and put $k=\wt(u)$.

\begin{lemma}[exact pointed reduction]\label{lem:phase}
On the rotation class of $u$, the identification $j\leftrightarrow R^ju$
conjugates $T_n$ to
\begin{equation}\label{eq:phase}
 \phi_u(j)=
 \begin{cases}
 j+k,&u_j=1,\\
 j-k,&u_j=0,
 \end{cases}
 \qquad j\in\Z/d\Z.
\end{equation}
\end{lemma}

\begin{proof}
The weight is constant on the rotation class.  At $R^ju$, the first bit is
$u_j$.  Equations \eqref{eq:branches} therefore give increments $k$ and
$n-k$.  Since $d\mid n$, the latter is $-k$ modulo $d$.
\end{proof}

Put
\[
 h=\gcd(k,d),\qquad L=d/h.
\]
The undirected Cayley graph on $\Z/d\Z$ with generator $k$ consists of $h$
cycles of length $L$.  For a component with representative $r$, order its
vertices as $r,r+k,\ldots,r+(L-1)k$ and set
$b_q=u_{r+qk}$, with $q\in\Z/L\Z$.  Thus a bit $1$ points forward and a bit
$0$ points backward in this generator order.

\begin{theorem}[complete generator-component dynamics]\label{thm:component}
Each generator component is classified as follows.
\begin{enumerate}[(i)]
\item If $L=1$, its vertex is fixed.
\item If $L=2$, its two vertices form one directed two-cycle.
\item If $L\geq3$ and $(b_q)$ is constant, the component is one directed
      $L$-cycle.
\item If $L\geq3$ and $(b_q)$ is nonconstant, its recurrent components are
      exactly the cyclic occurrences $b_qb_{q+1}=10$, one two-cycle per
      occurrence.  Every other vertex is transient.  More precisely, its
      preperiod is
\begin{equation}\label{eq:pointdepth}
 \tau(q)=
 \begin{cases}
 \min\{t\geq0:(b_{q+t},b_{q+t+1})=(1,0)\},&b_q=1,\\
 \min\{t\geq0:(b_{q-t-1},b_{q-t})=(1,0)\},&b_q=0.
 \end{cases}
\end{equation}
Consequently the maximum preperiod in the component is one less than the
longest cyclic constant run in $(b_q)$.
\end{enumerate}
\end{theorem}

\begin{proof}
The first two cases follow because $+1$ and $-1$ in generator coordinates
are equal when $L\leq2$.  For $L\geq3$, a constant word points coherently
around the cycle and gives an $L$-cycle.

Now suppose the component word is nonconstant.  Whenever $b_qb_{q+1}=10$,
the vertex $q$ points to $q+1$ and $q+1$ points back to $q$, so this edge is
a directed two-cycle.  Starting from a $1$, successive arrows move forward
through its $1$-run until the terminal $1$ of such an edge is reached.
Starting from a $0$, they move backward through its $0$-run until the
initial $0$ of such an edge is reached.  This proves \eqref{eq:pointdepth}
and shows that every vertex reaches one of the listed edges.

No further directed cycle exists: a cycle that ever reverses direction
uses an edge in both directions and is one of the listed two-cycles, while a
cycle that never reverses would orient the whole component coherently,
contrary to nonconstancy.  On a constant run of length $s$, the farthest
vertex is $s-1$ arrows from its boundary two-cycle.  Taking the longest run
proves the final statement.
\end{proof}

The theorem is pointwise, not merely a period count: factor the pointed
necklace into its $h$ generator cycles, read their constant runs, and obtain
every tail and recurrent component without iterating \eqref{eq:literal}.
Several recurrent components can occur in one necklace.  For example, the
six rotations of $111000$ split into three directed two-cycles.

\section{Period inventory and sharp clock}

\begin{theorem}[all possible periods]\label{thm:periods}
For $n=1$, the possible-period set is $\{1\}$.  For every $n\geq2$, it is
\begin{equation}\label{eq:periods}
 \{1,2\}\ \cup\ \{\ell:\ell\mid n,\ 3\leq\ell<n\}.
\end{equation}
\end{theorem}

\begin{proof}
Theorem~\ref{thm:component} shows that a period longer than two must be the
length $L=d/\gcd(k,d)$ of a constant generator component.  Hence $L\mid n$.
If $L=n$, then $d=n$ and $\gcd(k,n)=1$, so there is only one generator
component.  Its constancy would make the whole word constant, contradicting
least period $d=n>1$.  Thus every long period is a proper divisor of $n$.

Constant words realize period one.  For $n\geq2$, the rotation class with a
single $1$ contains a directed two-cycle, realizing period two.  It remains
to realize a proper divisor $L\geq3$.  Write $n=gL$, where $g\geq2$, and
take the $n$-word of weight $g$ whose $1$-support is
\begin{equation}\label{eq:witnesssupport}
 \{0,1,\ldots,g-2,g\}\subset\Z/n\Z.
\end{equation}
The cyclic zero gap after position $g$ is uniquely longest, so a rotation
fixing the word must fix that gap and hence is trivial; the word has least
period $n$.  Under generator $k=g$, the component congruent to $g-1$ modulo
$g$ is all zero.  It is therefore a directed cycle of length $n/g=L$ by
Theorem~\ref{thm:component}.  This realizes every period in
\eqref{eq:periods}.
\end{proof}

Let $H_n$ be the maximum preperiod of $T_n$ over $\F^n$.

\begin{theorem}[sharp clock and deepest census]\label{thm:clock}
One has
\begin{equation}\label{eq:clock}
 H_1=0,\qquad H_n=n-2\quad(n\geq2).
\end{equation}
For every $n\geq3$, exactly two states have preperiod $n-2$, namely
\begin{equation}\label{eq:deepest}
 010^{n-2}\quad\text{and its bitwise complement}.
\end{equation}
At the boundaries, both states are deepest for $n=1$ and all four states
are deepest for $n=2$.
\end{theorem}

\begin{proof}
In a nonconstant generator component of length $L$, a constant run has
length at most $L-1$.  Theorem~\ref{thm:component} bounds its preperiod by
$L-2\leq n-2$; constant components contain no transient vertex.  This gives
the upper bound for $n\geq2$.  For $n\geq3$, the first word in
\eqref{eq:deepest} has weight one and, in generator order, a cyclic zero run
of length $n-1$.  Its displayed pointer is $n-2$ steps from the recurrent
$10$ edge.  Complementation commutes with \eqref{eq:literal}, so the second
word is another witness.

Equality in the upper bound forces $L=n$ and a constant run of length
$n-1$.  Hence the word has exactly one minority bit.  There are precisely
two such rotation necklaces, one for each minority symbol, and in each
there is exactly one pointer farthest from the recurrent $10$ edge.  These
are the two words in \eqref{eq:deepest}.  Directly, both length-one words are
fixed, while at length two the constants are fixed and the two nonconstant
words form a two-cycle.  The boundary claims follow.
\end{proof}

The numerical value $n-2$ in \eqref{eq:clock} is not credited: P166 already
has that sharp value.  The retained statement is the generator-run proof
together with the exact two-state last shell.

\section{Every-target fibres and fixed census}

For a target $y\in\F^n$, subscripts such as $y_{-k}$ are interpreted modulo
$n$.  Let $N_r(n,k)$ be the number of weight-$k$ targets having exactly
$r$ one-step predecessors.

\begin{theorem}[labelled inverse atlas]\label{thm:fibres}
If $y$ is constant, then $T_n^{-1}(y)=\{y\}$.  If $0<k=\wt(y)<n$, then its
complete predecessor list is
\begin{align}
 x_1&=R^{-k}y &&\text{if and only if }y_{-k}=1,\label{eq:inverseone}\\
 x_0&=R^{k}y  &&\text{if and only if }y_k=0.\label{eq:inversezero}
\end{align}
Here the subscripts label the first bit of the source, so present sources
are distinct.  In particular every fibre has size $0$, $1$, or $2$.

For $0<k<n$, if $n\mid2k$, then
\begin{equation}\label{eq:halfweight}
 N_1(n,k)=\binom nk,\qquad N_0(n,k)=N_2(n,k)=0.
\end{equation}
Otherwise,
\begin{equation}\label{eq:fibrehist}
 N_0(n,k)=N_2(n,k)=\binom{n-2}{k-1},\qquad
 N_1(n,k)=\binom nk-2\binom{n-2}{k-1}.
\end{equation}
The two constant targets contribute two more to the global one-fibre
count, and
\begin{equation}\label{eq:image}
 |\Imap T_n|=2+\sum_{k=1}^{n-1}
 \left[\binom nk-
 \mathbf1_{\{n\nmid2k\}}\binom{n-2}{k-1}\right].
\end{equation}
\end{theorem}

\begin{proof}
A source beginning in $1$ must use the first branch of
\eqref{eq:branches}; undoing it gives the sole candidate $R^{-k}y$, whose
first bit is $y_{-k}$.  Similarly, a source beginning in $0$ uses rotation
$n-k$, whose inverse is $R^k$; its first bit is $y_k$.  This proves
\eqref{eq:inverseone}--\eqref{eq:inversezero}.  The two candidates, when
present, start with different bits.

The inspected positions $-k$ and $k$ coincide precisely when $n\mid2k$.
Then exactly one of the complementary branch conditions holds, proving
\eqref{eq:halfweight}.  Otherwise they are distinct.  Fibre zero prescribes
the ordered pair $(y_{-k},y_k)=(0,1)$, and fibre two prescribes $(1,0)$.
In either case the remaining $n-2$ positions contain $k-1$ ones, giving
$\binom{n-2}{k-1}$.  Subtraction gives the one-fibre count.  Finally,
subtracting the zero-fibre targets from each weight layer and adding the two
constants proves \eqref{eq:image}.
\end{proof}

Fixed-density necklace and primitive-word enumeration are classical
\cite{GrosekHromada2016,HeckenbergerSawada2018,Mestrovic2018}.  We use that
background only to record a supporting fixed census.

\begin{theorem}[M\"obius fixed census]\label{thm:fixed}
Let $\mu$ be the M\"obius function and define
\begin{equation}\label{eq:primitive}
 A(d,j)=\sum_{e\mid\gcd(d,j)}\mu(e)
          \binom{d/e}{j/e}.
\end{equation}
Then the number of fixed points of $T_n$ is
\begin{equation}\label{eq:fixed}
 |\Fix(T_n)|=
 \sum_{d\mid n}\ \sum_{j=0}^{d}
 \mathbf1_{\{d\mid(n/d)j\}}A(d,j).
\end{equation}
\end{theorem}

\begin{proof}
By M\"obius inversion, $A(d,j)$ counts length-$d$ binary linear words of
least period $d$ and weight $j$.  Repeating such a primitive block $n/d$
times gives a length-$n$ word of least period $d$ and total weight
$k=(n/d)j$, and every length-$n$ word arises uniquely in this way.

If the pointed first bit is $1$, the word is fixed exactly when $d\mid k$.
If it is $0$, it is fixed exactly when $d\mid n-k$, which is equivalent
because $d\mid n$.  Thus the fixed condition is independent of the pointer
and is exactly the indicator in \eqref{eq:fixed}.  Summing the primitive
blocks proves the formula.
\end{proof}

Primitive-word enumeration and M\"obius inversion in
\eqref{eq:primitive}--\eqref{eq:fixed} receive zero contribution credit.

\section{Exact control and lifecycle}

A paper-local, standard-library author/scout-derived regression control
starts from the literal update
\eqref{eq:literal} and enumerates all $2^n$ words for every $1\leq n\leq18$.
It checks carrier closure and weight preservation; every target in
\eqref{eq:inverseone}--\eqref{eq:inversezero}; the full fibre histogram;
every pointed-necklace conjugacy; every component tail and recurrent cycle;
the possible-period set; the sharp clock and deepest census; and
\eqref{eq:fixed}.  The frozen transcript reports $2{,}828{,}503$ assertions
and has SHA-256
\begin{center}
\small\ttfamily
3d0947a4df32f8e583e28d1964a52523602d61c64dde7b259bfdd15e71e4003b.
\end{center}
Finite enumeration is a falsification and regression control, not a
substitute for any all-parameter proof.

The retained results are precisely the
$\mathord{\pm}k$ pointed-necklace component theorem, the period inventory,
the sharp clock with deepest two, the every-target $0/1/2$ fibres, and the
M\"obius fixed census.  No external action is authorized.  The terminal
status is
\[
 \boxed{\texttt{AMBER\_INTERNAL\_NEAR\_P166 / HOLD\_EXTERNAL}}.
\]

\bibliographystyle{plainnat}
\bibliography{references}

\end{document}
